% исправления 27.05.2005
\chapter*{Предисловие}
\addcontentsline{toc}{schapter}{Предисловие}
\markboth{Предисловие}{Предисловие}

\begin{flushright}
\textit{Нашему учителю,}\\[1ex]
\textit{ВЛАДИМИРУ АНДРЕЕВИЧУ УСПЕНСКОМУ}\\[2ex]
%\textit{к 70-летию}\\
\end{flushright}

Предлагаемая вашему вниманию книга написана по материалам лекций
для младшекурсников, которые читались авторами в разные годы на
механико\д математическом факультете МГУ. (В эту серию также
входят книги \лк Начала теории множеств\пк\ и \лк Вычислимые
функции\пк.)

Центральная идея математической логики восходит ещё к Лейбницу и
состоит в том, чтобы записывать математические
утверждения в виде последовательностей символов и
оперировать с ними по формальным правилам. При этом правильность
рассуждений можно проверять механически, не вникая в их смысл.

Усилиями большого числа математиков и логиков второй
половины XIX и первой половины XX~века (Буль, Кантор, Фреге,
        \glossary{\Буль}%
        \glossary{\Кантор}%
        \glossary{\Фреге}%
        \glossary{\Пеано}%
        \glossary{\Рассел}%
        \glossary{\Уайтхед}%
        \glossary{\Цермело}%
        \glossary{\Гильберт}%
        \glossary{\фонНейман}%
        \glossary{\Гёдель}%
        \glossary{\Френкель}%
Пеано, Рассел, Уайтхед, Цермело, Френкель, Гильберт, фон Нейман,
Гёдель и другие) эта программа была в основном выполнена.
Принято считать, что всякое точно сформулированное
математическое ут\-вер\-ж\-де\-ние можно записать формулой теории
множеств (одной из наиболее общих формальных теорий), а всякое
строгое математическое доказательство преобразовать в формальный
вывод в этой теории (последовательность формул теории множеств,
подчиняющуюся некоторым простым правилам). В каком\д то смысле это
даже стало определением: математически строгим считается такое
рассуждение, которое можно перевести на язык теории множеств.

Так что же, теперь математики могут дружно уйти на пенсию,
поскольку можно открывать математические теоремы с помощью
компьютеров, запрограммированных в соответствии с формальными
правилами теории множеств? Конечно, нет, причём сразу по
нескольким причинам.

Начнём с того, что машина, выдающая с большой скоростью
математические теоремы (и их доказательства), хотя и возможна,
но бесполезна. Дело в том, что среди этих верных утверждений
почти все будут неинтересными. Формальная логика говорит, какие
правила надо соблюдать, чтобы получать верные результаты, но не
говорит, в каком порядке их надо применять, чтобы получить
что\д то интересное.

Казалось бы, мы можем запустить машину и ждать, пока она не
докажет интересующее нас утверждение (пропуская все остальные).
Проблема в том, что формальное доказательство сколько\д нибудь
содержательной теоремы настолько длинно, что прочесть его
человек не в состоянии. Представьте себе доказательство, которое
состоит из миллионов формально правильных шагов, в котором мы
можем проверить каждый отдельный шаг, но так и не понимаем, что
происходит\т много ли в нём~проку?

На самом деле прок всё\д таки есть: мы узнаём, что
доказываемое утверждение верно, хотя так и не понимаем, почему.
Так что и такая машина была бы полезна. Увы, и этого сделать не
удаётся, поскольку на поиск доказательства сколько\д нибудь
сложного утверждения известными сейчас методами требуется
астрономически большое время (даже если представить себе, что
машина работает с предельно возможной по законам физики
скоростью).

Можно умерить амбиции и поставить  задачу попроще: пусть
машина проверяет доказательства, записанные человеком по
правилам формальной логики. Если машина не может помочь нам
что\д то открыть, пусть она хотя бы проверит, не пропустили ли мы
какого\д то шага рассуждения.

Из всех перечисленных задач эта выглядит наиболее реалистичной.
К сожалению, пока что работы и в этом направлении не ушли
далеко: формальная запись доказательства в виде, пригодном для
машинной проверки, является долгим и скучным делом, на которое у
большинства математиков не хватает энтузиазма и терпения. А разработать
удобные средства такой записи пока не удалось.

Короче говоря, революционная программа Лейбница построения
формальных оснований математики осуществилась, но незаметно: под
здание математики подвели новый (и довольно прочный) фундамент,
но большинство жильцов про это до сих пор не знают.

Так что же, математическая логика бесполезна? Ни в коем
случае: она не только удовлетворяет естественный философский
интерес к основаниям математики, но и содержит множество
красивых результатов, которые важны не только для математики,
но и для computer science.

В этой книжке мы расскажем об одном из центральных понятий
математической логики\т языках и исчислениях первого порядка. В
этих языках используются логические связки \лк и\пк, \лк или\пк,
\лк если\dots\  то\dots\пк, а также кванторы \лк для всех\пк\ и
\лк существует\пк. Оказывается, что этих средств достаточно для
формализации математических теорий и что можно построить простые
формальные правила, полностью отражающие смысл этих логических
средств.

Авторы пользуются случаем ещё раз поблагодарить своего учителя,
Владимира Андреевича Успенского, лекции, тексты и высказывания
которого повлияли на них (и на содержание этой книги), вероятно,
даже в большей степени, чем авторы это осознают.

При подготовке текста использованы записи
А.\,Ев\-фимь\-евс\-ко\-го и А.\,Ро\-ма\-щен\-ко (который также
прочёл предварительный вариант книги и нашёл там немало ошибок).

Оригинал\д макет книги был подготовлен
%её редактором,
В.\,В.\,Шу\-ва\-ло\-вым; без его настойчивости (вплоть до
готовности разделить ответственность за ошибки) оригинал\д макет
вряд ли появился бы к какому\д либо сроку. Он же вместе с
М.\,А.\,Уша\-ко\-вым (нашедшим несколько существенных ошибок)
подготовил предметный указатель. Мы признательны также
К.\,С.\,Ма\-ка\-ры\-че\-ву и Ю.\,С.\,Ма\-ка\-ры\-че\-ву, которые внимательно
прочли вёрстку книги и нашли там немало опечаток.

Авторы признательны \'Ecole Normale Sup\'erieure de Lyon
(Франция) за поддержку и гостеприимство во время написания этой
книги.

Первое издание книги стало возможным благодаря Российскому фонду
фундаментальных исследований, а также И.\,В.\,Ященко, который
уговорил авторов подать туда заявку.

Наконец, мы благодарим сотрудников, аспирантов и студентов
кафедры математической логики мехмата МГУ (особая
благодарность\т М.\,Р.\,Пентусу, указавшему два десятка
опечаток), а также всех участников наших лекций и семинаров и
читателей предварительных вариантов этой книги.

В третьем издании добавлены формулировка и доказательство
теоремы Чёрча о неразрешимости исчисления предикатов (по ошибке
отсутствовашие в предыдущих изданиях), а также дополнена
информация в именном указателе. В четвёртом издании, помимо
изменения формата вёрстки и использования шрифтов LH, исправлены
некоторые опечатки (на которые нам любезно указали Н.~Маслов, А.~Гусаков, В.~Патков).

Просим сообщать о всех ошибках и опечатках авторам (электронные
адреса \texttt{ver} at \texttt{mccme} dot \texttt{ru},
\texttt{nikolay} dot \texttt{vereshchagin} at \texttt{gmail} dot
\texttt{com}; \text{sasha} dot \texttt{shen} at \texttt{gmail} dot \texttt{com},
\texttt{alexander} dot \texttt{shen} at \texttt{lirmm} dot \texttt{fr}; почтовый адрес: Москва,
119002, Большой Власьевский пер., 11, Московский центр
непрерывного математического образования).

\begin{flushright}
        %
\textit{Н.\,К.\,Верещагин, А.\,Шень}
        %
\end{flushright}
